Il reticolo olografico è un tipo di reticolo di difrazione e quindi è una successione di "fenditure" parallele e rettilinee. Il termine olografico nasce dal fatto che la tecnica con cui viene realizzata la griglia di difrazione consiste nel avere una pellicola trasparente in cui viene impressa una figura di interferenza.

Quando una luce non monocromatica colpisce un reticolo di difrazione ad una distanza $z$, la luce uscente non rimane inalterata ma viene scomposta nelle sue componenti cromatiche. Ogni punto sul reticolo viene colpito dalla luce che viene dalla sorgente bianca con un angolo diverso. Da ogni punto quindi verrà difratta la luce con angoli diversi, come è mostrato nell'immagine. Ma il risultato è che questi raggi divergono come se provenissero da un punto posto nel piano della sorgente ma spostato di $\Delta x$ e questo accade per tutti i colori. La distanza tra la sorgente virtuale e la sorgente reale è legata dalla seguente equazione:
\begin{equation}
    \Delta x= z\tan{\phi}\approx z\sin{\phi}=\frac{z\lambda}{d}
\end{equation}
Dove z è la distanza tra la sorgente e il reticolo, infatti più il reticolo è lontana dalla sorgente più l'angolo che separa la sorgente virtuale e la sorgente reale è maggiore, $\lambda$ è la lunghezza d'onda della sorgente virtuale ed infine $d$ è il passo del reticolo.

Durante l'esperiemento bisogna misurare la lunghezza d'onda attraverso il reticolo olografico con il passo $d$ noto con la seguente equazione:
\begin{equation}
    \lambda=\frac{\Delta x d}{z}
\end{equation}
E successivamente, utilizzando la stessa configurazione, sostituire il reticolo olografico con il secondo e ricavere $d'$:
\begin{equation}
    d'=\frac{z\lambda}{\Delta x}
\end{equation}